\documentclass[12pt,a4paper]{article}

% -------------------------------------------------
% PACKAGES
% -------------------------------------------------
\usepackage[a4paper, margin=1in]{geometry}
\usepackage{setspace}
\usepackage{enumitem}
\usepackage{titlesec}
\usepackage{array}

% -------------------------------------------------
% DOCUMENT SETTINGS
% -------------------------------------------------
\setlength{\parindent}{0pt}
\setlength{\parskip}{6pt}

\titleformat{\section}
  {\normalfont\bfseries\large}
  {}
  {0pt}
  {}

\begin{document}

% -------------------------------------------------
% TITLE
% -------------------------------------------------

\begin{center}
    {\Large \textbf{PROJECT JOURNAL }}\\[10pt]

    {\large \textbf{BlurGuard -- Real-Time Screen-Share Sensitive Data Censor}}
\end{center}

\vspace{10pt}

% -------------------------------------------------
% JOURNAL DETAILS
% -------------------------------------------------

\begin{tabular}{@{}p{0.48\textwidth}p{0.48\textwidth}@{}}
\textbf{Journal :} 03 &
\raggedleft \textbf{} September 2026
\end{tabular}

\vspace{12pt}

% -------------------------------------------------
% TEAM MEMBERS
% -------------------------------------------------

\textbf{Team Members:}

\begin{itemize}[leftmargin=1.5cm]
    \item Lakshita
    \item Rhythm
    \item Yashjot
\end{itemize}

\vspace{5pt}

% -------------------------------------------------
% WORK COMPLETED
% -------------------------------------------------

\textbf{\large Work Completed}

\vspace{4pt}
This week we worked on the next step which was documentation . I created a small plan consisting of list of tasks we will work on:
\par
[ Plan : all diagrams(NO AI, All diagrams must have same modules and user list that should be fixed beforehand):\par
1. Use case diagram [rhythm]\par
2. Use-case templates for each module [lakshita]\par
3. Activity diagram [rhythm]\par
4. Swimlane diagram [lakshita]\par
5. Sequence diagram[lakshita]\par
6. Collaboration diagram [lakshita]\par
+ interface with some modules working [rhythm, yash, lakshita]
 + report [yash] ]

Then we divided the tasks as specified in above plan. I divided the project into 12 modules and mentioned what all functioning will be done in them which actors are accessing them.
Then I worked with use case diagram, activity diagram, finalizing and uploading journals and started with development part.

\vspace{8pt}

% -------------------------------------------------
% OUTCOME
% -------------------------------------------------

\textbf{\large Outcome}

\vspace{4pt}

Started Documentation and implementation procedures, where we finalised modules and worked with uml diagrams and interface with basic modules working .

\vspace{8pt}

% -------------------------------------------------
% NEXT PLANNED WORK
% -------------------------------------------------

\textbf{\large Next Planned Work}

\vspace{4pt}

To create DFD (level 0,level 1,level 2) and add functionality of more modules and make it work for real time data detection and masking on application level. 

\vfill

% -------------------------------------------------
% FOOTER
% -------------------------------------------------

\begin{flushright}
    \textit{Blurguard Project Journal}
\end{flushright}

\end{document}